\chapter*{Declaración de uso responsable de inteligencia artificial}
\addcontentsline{toc}{chapter}{Declaración de uso responsable de inteligencia artificial}

Durante la elaboración de este Trabajo Fin de Grado se ha recurrido a herramientas de
\textbf{inteligencia artificial (IA)} generativa como apoyo en dos ámbitos diferenciados:
el desarrollo del software CloudPorter y la redacción de la presente memoria. A continuación
se detallan las herramientas utilizadas, los usos concretos realizados y los mecanismos de
validación empleados para garantizar la corrección y coherencia del trabajo entregado.

\section*{Herramientas utilizadas}

La única herramienta de IA empleada a lo largo del proyecto ha sido \textbf{Claude Code}
\cite{claudeCode}, un agente de programación e interacción en lenguaje natural desarrollado
por Anthropic, que opera desde la línea de comandos y tiene acceso al repositorio y al
entorno de desarrollo local.

\section*{Usos realizados}

El uso de Claude Code se divide en dos vertientes:

\subsection*{Apoyo al desarrollo de CloudPorter}

Claude Code se empleó como asistente de programación durante el ciclo de desarrollo iterativo
descrito en el capítulo~\ref{cap:desarrollo}. Los usos concretos incluyeron:

\begin{itemize}
    \item \textbf{Generación y revisión de código:} implementación de módulos del traductor,
    plantillas Jinja2 y la lógica del mapper determinístico, revisados y validados por el
    autor antes de cada integración.
    \item \textbf{Skills de flujo de trabajo:} se desarrollaron dos skills propias para
    Claude Code —\texttt{create-issue} y \texttt{github-workflow}— que encapsulan las
    convenciones del proyecto y asisten en la creación de issues y la revisión de workflows
    de GitHub Actions.
    \item \textbf{Reglas de contexto:} se configuró un fichero de reglas específico para
    el módulo de traducción (\texttt{.claude/rules/translator.md}) que documenta las
    convenciones de estructura del traductor y orienta al agente cuando trabaja en esa
    parte del código.
    \item \textbf{Apoyo en depuración y diseño:} consultas sobre decisiones de arquitectura,
    resolución de errores en la suite de tests y revisión de pull requests.
\end{itemize}

En todos los casos, el código generado fue revisado, adaptado y validado por el autor
mediante la suite de tests automatizados del proyecto —con un umbral mínimo de cobertura
del 80\%— y mediante la ejecución manual de los flujos de despliegue sobre AWS y Azure.

\subsection*{Apoyo a la redacción de la memoria}

Claude Code se utilizó como asistente de escritura académica para la redacción y revisión
de los capítulos de esta memoria. Los usos concretos incluyeron:

\begin{itemize}
    \item \textbf{Redacción de secciones:} generación de borradores de capítulos a partir
    de indicaciones del autor, respetando el enfoque narrativo del TFG y las convenciones
    de estilo definidas.
    \item \textbf{Revisión lingüística y de estilo:} corrección de coherencia, fluidez y
    adecuación al registro académico en castellano.
    \item \textbf{Skill de redacción:} se configuró una skill propia (\texttt{escritura-tfg})
    que encapsula las convenciones de voz, estructura y estilo propias de esta memoria,
    asegurando consistencia entre sesiones de trabajo.
    \item \textbf{Contexto persistente:} se mantiene un fichero \texttt{CLAUDE.md} con el
    enfoque narrativo del TFG, la estructura de capítulos y las instrucciones de redacción,
    que se carga automáticamente en cada sesión de trabajo con el agente.
\end{itemize}

Todo el contenido generado ha sido revisado críticamente por el autor, que ha verificado
la corrección técnica de las afirmaciones, la coherencia con el resto del documento y
la adecuación a los requisitos académicos. Las decisiones de contenido —qué incluir,
qué omitir, cómo enmarcar cada argumento— son en todo momento responsabilidad del autor.

\section*{Validación}

Los mecanismos de validación empleados para garantizar la calidad del trabajo son los
siguientes:

\begin{itemize}
    \item \textbf{Código:} suite de tests automatizados con cobertura mínima del 80\%,
    comprobaciones de tipos con mypy, análisis estático con Ruff y validación funcional
    mediante despliegue real en AWS y Azure.
    \item \textbf{Memoria:} revisión manual de cada sección generada, contraste de datos
    técnicos con la documentación oficial de las herramientas y fuentes bibliográficas
    citadas, y compilación continua del documento LaTeX para verificar su integridad formal.
\end{itemize}

\section*{Periodo de uso}

Las herramientas de IA descritas se emplearon durante todo el período de desarrollo e
implementación del proyecto, comprendido entre febrero y mayo de 2026.
